\chapter{Signal control and effective service}
\label{chap:signal-control}

\section{Controllers, phases and groups}
A controller defines cycle time, offset and ordered phases.  Each phase contains one or
more active signal groups.  Movements reference signal groups rather than phases
directly.  This permits one signal group to be active in more than one phase.

\section{Displayed green and intergreen}
For signal group $j$, displayed signal green is the duration for which the group is
active.  Intergreen belongs to transition/clearance time and is not automatically
counted as signal green.

For a simple non-overlapping phase sequence,
\begin{equation}
C=\sum_p(G_p+I_p),
\end{equation}
where $C$ is cycle time, $G_p$ is displayed phase green and $I_p$ is the intergreen
after phase $p$.

\section{Continuous signal-group runs and phase overlap}
If a signal group is active in consecutive phases, PySig treats the group as
continuously active across the internal phase boundary.  For example:

\begin{verbatim}
Phase P_A: SG_A, SG_B
Phase P_B:       SG_B, SG_C
\end{verbatim}

SG_B has one continuous signal run.  The intergreen after P_A is not inserted into
SG_B's service simply because the controller changes phase.  Ending displacement and
clearance treatment apply when the continuous SG_B run actually ends.

\section{Start displacement and end displacement}
PySig distinguishes displayed signal green $G$ from engineering effective green $g$.
With start displacement $d_s$ and end displacement $d_e$,
\begin{equation}
g=G-d_s+d_e,
\end{equation}
bounded to the interval $0\leq g\leq C$.

For 55 s of displayed green, 3 s start displacement and 3 s end displacement,
\begin{equation}
g=55-3+3=55\text{ s}.
\end{equation}
The duration is unchanged, but the effective service interval is displaced in time.
An asymmetric 3 s start displacement and 2 s end displacement gives 54 s of effective
green.

\section{Capacity implication}
The nominal fixed-time capacity is
\begin{equation}
c_n=s\frac{g}{C},
\end{equation}
where $s$ is lane saturation flow.  The corresponding nominal degree of saturation is
\begin{equation}
x_n=\frac{q}{c_n}=\frac{qC}{sg}.
\end{equation}
These relationships are consistent with conventional signal-capacity definitions.
Operational accessible capacity may be lower after lane interactions are considered.

\section{Priority and yielding movements}
A yielding movement may identify one or more priority movements and an active capacity
factor.  When no priority movement is active, the factor is 1.0.  While a configured
priority stream is effectively active, available yielding service is reduced by the
configured factor.  The factor is applied within the shared physical-lane supply; it
does not create separate full lane capacity for each movement.

The current factor formulation is an explicit planning approximation.  It is not a
full gap-acceptance model involving critical gap, follow-up headway, pedestrian
conflicts or opposing platoon structure.
